\gddchapter{Reflection analysis}

The role of this chapter is to provide reflections post interview and to refine the process for the next interview.

\section{Methodology reflecion}
\begin{enumerate}
    \item Wait time \& silence: 
    
    Did I provide the participant with time and space to voice out their opinions or did I fill the silence gap due to my own discomfort?

    Due to the fact that this was the first test I admit that I jumped in to fill the silence a bit too much. I take that as a reflection for future tests.
    \item Lexical mirroring 
    
    Identify 1-2 times where I used participants own words to deepen his narrative and answer.
\end{enumerate}

\section{Defensive self-probing}

\begin{QandA}
   \item Wait time \& silence: Did I provide the participant with time and space to voice out their opinions or did I fill the silence gap due to my own discomfort?
         \begin{answered}
            Due to the fact that this was the first test I admit that I jumped in to fill the silence a bit too much. I take that as a reflection for future tests.
         \end{answered}

   \item Lexical mirroring: Identify 1-2 times where I used participants own words to deepen his narrative and answer.

         \begin{answered}
         The participant used the word "zanurzenie" - it's a specific word that can mean immersion ``No to jak takie zanurzenie się w świecie'' (This is the immersion in the game world) and the researcher used the same word asking a qustion ``To znaczy, mówiłaś, że czułaś większe zanurzenie, tak?'' (This means that you felt a deeper immersion right?)
         \end{answered}

    \item The "Vessel trap": Did I treat the participant as a way to prove my internal bias that this system might be more immersive than traditional input?

         \begin{answered}
        I don't think I treated the participant as a way to prove my internal bias here. I thought that since the system was breaking so much that there would be no way that it is more immersive in the first place. The amount of embarasment for every single bug that the participant had to experience made me not focused on proving my internal bias, instead I tried to distance myself from my own creation mostly due to that embarasment that I felt.
         \end{answered}
\end{QandA}

\section{Sticky moments analysis}
Analysis of the meaning of the sticky moments noted in the chapter 3


\begin{enumerate}
    \item \textbf{Location description length}
    
    The player complained about the length of the location descriptions.
    The descripions are too long I agree. They should have been maybe 50\% of their length.


    \item \textbf{Item usage}
    
    It was not directly obvious how to use items in the keyboard version.

    The UX of that interaction can use some improvement. As is currently the item selection arrow is too small.


    \item \textbf{Reading time}
    
    The player returned to this: \textit{``I could be not reading this at all. This takes too long.''}


    \item \textbf{Stair naming}
    
    The naming convention for the stairs was not intuitive.
    The stairs are called "Stairway going up" and "Stairway going down". Therefore if the player wants to go up the stairs they have to travel to "Stariway going down". It's unitntuitive. 

    \item \textbf{Recording UX}
    
    The UX of the recording mechanic was not obvious.

    The player pressed the R button, released it and then said their command.


    \item \textbf{Stair image bug}
    
    The 1F stairs displayed the same image as 0F. I defended the game more than I should have.


    \item \textbf{Lounge entry}
    
    Double UX fail entering the lounge. The word lounge was problematic to say for the participant due to their accent.


    \item \textbf{Pharmacy pickup}
    
    The leg stabiliser refused to be picked up on the first attempt but succeeded on the second. Seems like a voice recognition error.


\end{enumerate}



\section{Refinment and iterative tweaks}

\begin{QandA}
    \item Do any questions need changes or refinements based on my observations?
    \begin{answered}
        Yes, the first Grand tour question about describing the user journey has left the player confused and didn't provide any tangible benefit in my opinion. The player just restated what they did in the game.
    \end{answered}

    \item Does the testing procedure need any changes based on the observations from this session?
    \begin{answered}
        I feel like the fact that the player was initially exposed to commands in version A of the game made them choose their words differently in the version B. The order of those experiences seems to be influencing one another. I'll verify this with future tests
    \end{answered}
\end{QandA}